\documentclass[10pt]{article}

\usepackage[utf8]{inputenc}

\usepackage[T1]{fontenc}

\usepackage{amsmath}

\usepackage{amsfonts}

\usepackage{amssymb}

\usepackage[version=4]{mhchem}

\usepackage{stmaryrd}

\usepackage{graphicx}

\usepackage[export]{adjustbox}

\graphicspath{ {./images/} }

\usepackage{mathrsfs}



\begin{document}

образом определить подклассы класса $H^{*}$, в к-рых разыскиваются решения данного и союзного с ним C. и. у., то остаются в силе и теоремы Нётера (см. [1]). Указанные выше результаты обобщены в различных направлениях. Показано (см. [1]), что при нек-рых условиях они остаются справедливыми и в случае



\includegraphics[max width=\textwidth, center]{2023_07_08_66f6f0101c99a26ca1c5g-1}

конечного числа гладких разомкнутых дуг, к-рые могут попарно пересекаться только по своим концам). С. и. у. исследованы также в функциональных пространствах Лебега $L_{p}(\Gamma)$ и $L_{p}(\Gamma, \rho)$, где $p>1$, а $\rho-$ нек-рый вес (см. [4] - [7]). В работах [4]-[6] приведены результаты, к-рые непосредственно обобщают выше сформулированные утверждения.



Пусть уравнение простой спрямляемой линии $\Gamma$ есrь $t=t(s), 0 \leqslant s \leqslant \gamma$, где $s$ - дуга этой линии, отсчитываемая от нек-рой фиксированной точки на ней, $\gamma$ - длина $Г$. Говорят, что функция $f$, определенная на $\Gamma$, почти всюду конечна, измерима, интегрируема и т. д., если функция $f(t(s))$ на сегменте $[0, \gamma]$ обладает соответствующим свойством. Интеграл Лебега от $f$ на $\Gamma$ ошределяют равенством



\[

\int_{\Gamma} f(t) d t=\int_{0}^{\nu} f(t(s)) t^{\prime}(s) d s .

\]



Через $L_{p}(\Gamma)$ обозначается множество измеримых на $\Gamma$ функций таких, что $|f|^{p}$ интегрируема на $\Gamma$. Функциональный класс $L_{p}(\Gamma), p \geqslant 1$, шреврашается в банахово пространство, если норму элемента $f$ определить равенством



\[

\|f\|=\left(\int_{\Gamma}|f|^{p} d s\right)^{1 / p}

\]



Если в уравнениях (3), (7), в к-рых равенства соблюдаются почти всюду, коэффициенты $a, b$ непрерывны и удовлетворяют условию нормальности, $f, g \in L_{p}(\Gamma)$, $p>1$, то остаются в силе утверждения 1) и 2), если в них класс $H$ заменить классом $L_{p}(\Gamma), p>1$. Далее, если решения уравнения $K \varphi=f$, где оператор $K$ имеет вид (2), разыскиваются в банаховом пространстве $L_{p}(\Gamma), p>1$, а решения союзного однородного с ним уравнения $K^{\prime} \psi=0-$ в пространстве $L_{p^{\prime}}(\Gamma)$, где $p^{\prime}=$ $=p /(p-1)$, то остаотся в силе и теоремы Нётера, причем $V$ может быть любым вполне непрерывным оператором в $L_{p}(\Gamma)$.



Когда Г является конечной совокупностью разомкнутых линий или $\Gamma$ замкнута, но коэффициенты С. и. у. терпят разрывы, то решения уравнений часто разыскиваются в весовых функциональных пространствах $L_{p}(\Gamma, \rho), p>1 \quad\left(f \in L_{p}(\Gamma, \rho) \Longleftrightarrow \rho f \in L_{p}(\Gamma)\right)$. При определенных условиях относительно весовой функции $\rho$ остаются сіраведливыми результаты, аналогичные вышеприведенным.



Задача регуляризации. Одной из основных задач, возникаюгцих при построении теории С. и. у., является задача регуляризации, т. е. задача приведения C. и. у. к уравнению Фредгольма.



Пусть $E$ и $E_{1}$ - банаховы пространства, к-рые могут и совпадать, $A$ - линейный ограниченный оператор $A: E \rightarrow E_{1}$. Ограниченный оператор В наз. л е в ы м ре г уля р и затором опе ратора $A$, если $B A=I+V$, где $I, V-$ тождественный и вполне непрерывный операторы в $E$. Если уравнения $A \varphi=f$ и $B A \varphi=$ $=B f$ эквивалентвы, каков бы ни был элемент $f \in E_{1}$, то $B$ наз. л е вым эк в и в ал ле ты ы м ре г уля я и з а тором о пе р а то ра $A$. Ограниченный оператор $B$ наз. п р а в ы м $\mathrm{p}$ е г у л я р и зато ро м о пе р а то р а $A$, если $A B=I_{1}+V_{1}$, где $I_{1}, V_{1}$ - тождественный и вполне непрерывный операторы в $E_{1}$ соответственно. Если при любом $f \in E_{1}$ уравнения $A \varphi=f$ и $A B \psi=f$ одновременно разрешимы или неразрешимы, причем в случае разрешимости между их решениями существует связь $\varphi=B \psi$, то $B$ на3. п р а вы м әк ви в а ле н т ны м р е гуляри а т т ром о пе р а то р а $A$. Если $B$ является одновременно левым и правым регуляризатором оператора $A$, то его наз. д в у с т о р о н н и м р е г уля я и з а торо м, или просто ре г уля р и з ат о р о м, о п е р а т о р а $A$. Говорят, что оператор $A$ до пускае т регуля ри зацию ле вую, п р а в ую, д в у с то р о нню ю, әк в и в а ле нтн у ю, если существует его регуляризатор (соответственно левый, правый, двусторонний, әквивалентный).



Пусть $K$ - оператор, определяемый равенством (2), где $\Gamma$ - замкнутая простая гладкая линия, $a, b$ суть $H$-функции (или непрерывные функции), удовлетворяющие условию нормальности, $\mathscr{V}-$ вполне непрерывный оператор в пространстве $L_{p}(\Gamma), p>1$. Тогда в этом последнем пространстве оператор $K$ имеет бесчисленное множество регуляризаторов, среди к-рых находится, напр., оператор



\[

M=\frac{a}{a^{2}-b^{2}} I-\frac{b}{a^{2}-b^{2}} S .

\]



Для того чтобы оператор $K$ допускал левую әквивалентную регуляризацию, необходимо и достаточно, чтобы индекс $x$ оператора $K$ был неотрицательным [7]. В качестве эквивалентного левого регуляризатора можно взять оператор $M$. Если $x<0$, то оператор $K$ допускает правую эквивалентную регуляризацию, к-рую можно осуществить с помощью оператора $M$ (cм. [1]).



Системы С. и. у. Если в (1) $a, b, k$ - квадратные матрицы $n$-го порядка, рассматриваемые как матрицы линейных преобразований искомого вектора $\varphi=\left(\varphi_{1}\right.$, $\left.\ldots, \varphi_{n}\right)$, а $f=\left(f_{1}, \ldots, f_{n}\right)-$ известный вектор, то (1) наз. системой С. и. у. Система (1) наз. нормального тиіа, если матрицы $A=a+b$ и $B=a-b$ неособенные на $\Gamma$, то есть $\forall(t \in \Gamma) \operatorname{det} A \neq 0, \operatorname{det} B \neq 0$.



Теоремы Нётера остаются в силе для системы С. и. у. в классе $H$ (см. [1], [3]) и обобщаются на случай функциональных пространств Лебега (см. [4], [5]). В отличие от одного уравнения характеристич. системы C. и. у. в общем случае не решаются в квадратурах, но для вычисления индекса и в этом случаен айдена (см. [1]) формула, аналогичная формуле (9):



\[

\text { ind } K=\frac{1}{2 \pi}\left[\arg \operatorname{det} A^{-1} B\right]_{\Gamma} \text {. }

\]



В случае системы С. и. у. задачи регуляризации (см. [3]) аналогичны задачам регуляризации для одного C. и. $\mathrm{y}$.



В различных постановках исследованы как одно С. и. у., так и их системы, когда нарушаются условия нормальности (см. [11] и указанную там библиографию).



Многомерные С. и. у. Так называются уравнения вида $a(t) \varphi(t)+\int_{\Gamma} \frac{g(t, \vartheta)}{r^{m}} \varphi(\tau) d \tau+(V \varphi)(t)=f(t), t \in \Gamma$,



где $\Gamma$ - область евклидова пространства $E_{m}, m>1$; $\Gamma$ может быть конечной или бесконечной, в частности может совпасть с $E_{m} ; t, \tau$ - точки пространства $E_{m}$, $r=|t-\tau|, \quad \vartheta=(\tau-t) / r, \quad d \tau-$ әлемент объема в пространстве $E_{m}, V$ - вполне непрерывный оператор в банаховом функциональном пространстве, в к-ром ищется решение $\varphi ; a, g$ - заданные функции; несобственный сингулярный интеграл понимается в смысле главного значения, т. е.



\[

\int_{\Gamma} \varphi(\tau) \frac{g(t, \vartheta)}{r^{m}} d \tau=\lim _{\varepsilon \rightarrow 0} \int_{\Gamma \backslash\{r<\varepsilon\}} \varphi(\tau) \frac{g(t, \vartheta)}{r^{m}} d \tau

\]



точка $t$ наз. п о л ю с о м, функция $g(t, \vartheta)-\mathrm{x}$ а $\mathbf{p}$ а кт е р и с т и к о й, функция $\varphi$ - п л о т н о с т ь ю син-





\end{document}